\documentclass[a4paper,12pt]{article}
\usepackage[T2A]{fontenc}
\usepackage[utf8]{inputenc}
\usepackage[russian]{babel}
\usepackage{amsmath,amssymb,mathtools}
\usepackage{PTSans}
\renewcommand{\familydefault}{\sfdefault}
\usepackage{icomma}
\usepackage[a4paper,margin=20mm,headheight=25pt,headsep=7mm,footskip=12mm]{geometry}
\usepackage{microtype}
\usepackage{tikz}
\usetikzlibrary{arrows.meta,positioning,shapes.geometric,calc}
\usepackage{tabularx,booktabs}
\usepackage{enumitem}
\usepackage{xcolor}
\usepackage{parskip}
\usepackage{fancyhdr}
\usepackage{setspace}
\usepackage[unicode=true,colorlinks=true,linkcolor=black]{hyperref}

\definecolor{accent}{HTML}{147D7E}
\definecolor{darkaccent}{HTML}{0B5253}
\definecolor{ink}{HTML}{293538}
\definecolor{quiet}{HTML}{697578}
\definecolor{hairline}{HTML}{D5DDDD}
\color{ink}
\setstretch{1.12}
\setlength{\parindent}{0pt}
\setlength{\parskip}{5pt}
\setlist[itemize]{leftmargin=1.5em,itemsep=0.15em,topsep=0.25em}
\setlist[enumerate]{leftmargin=1.8em,itemsep=0.4em,topsep=0.30em}
\AtBeginDocument{\setlength{\abovedisplayskip}{7pt}\setlength{\belowdisplayskip}{7pt}}
\pagestyle{fancy}
\fancyhf{}
\fancyhead[L]{\small\color{quiet}Уравнения. Метод замены переменной}
\fancyhead[R]{\small\color{quiet}12.10.26}
\fancyfoot[C]{\small\color{quiet}\thepage}
\renewcommand{\headrulewidth}{0.4pt}
\renewcommand{\headrule}{\hbox to\headwidth{\color{hairline}\leaders\hrule height \headrulewidth\hfill}}
\newcommand{\key}[1]{{\color{darkaccent}\textbf{#1}}}
\newcommand{\note}[1]{{\small\color{quiet}#1}}
\newcommand{\headsec}[1]{\section*{\color{darkaccent}#1}\vspace{-5pt}{\color{hairline}\rule{\linewidth}{0.55pt}}\par\vspace{1pt}}
\newcommand{\headsub}[1]{\subsection*{\color{accent}#1}}

\tikzset{
  fstart/.style={ellipse,draw=darkaccent,line width=0.9pt,fill=white,
      minimum width=45mm,minimum height=11mm,align=center,font=\small},
  faction/.style={rounded corners=2.2pt,draw=accent,line width=0.85pt,fill=white,
      text width=104mm,minimum height=14mm,align=center,font=\small,inner sep=3.2mm},
  farrow/.style={-{Stealth[length=2.7mm]},draw=darkaccent,line width=0.85pt}
}

\begin{document}
\hypersetup{pageanchor=false}
\begin{titlepage}
\thispagestyle{empty}
\centering
\vspace*{28mm}
{\large\color{quiet}Чек-лист}\par
\vspace{7mm}
{\fontsize{29}{34}\selectfont\bfseries\color{darkaccent}
Решение уравнений\\[2mm]методом замены переменной}\par
\vspace{9mm}
{\Large\color{accent}Повторяющиеся выражения\\и обратная замена}\par
\vspace{17mm}
{\color{hairline}\rule{0.73\linewidth}{0.7pt}}\par
\vspace{10mm}
{\large Ярослав Верещагин}\par
\vspace{3mm}
{\normalsize 12.10.26}\par
\vfill
\begin{minipage}{0.82\linewidth}
\centering\normalsize
Научиться замечать скрытую структуру уравнения, выбирать удачную замену,
учитывать ограничения и находить все допустимые корни.
\end{minipage}
\vspace{13mm}
\begin{tikzpicture}
  \draw[accent!65,line width=1.4pt]
   (-6.5,0) .. controls (-3.5,0.7) and (-1.1,-0.45) .. (1.4,0.1)
     .. controls (3.1,0.53) and (4.7,0.18) .. (6.5,0.55);
\end{tikzpicture}\par
\vspace{9mm}
{\small\color{quiet}Лёвин Артём Александрович эксклюзивно для Ярослава Верещагина}
\end{titlepage}
\hypersetup{pageanchor=true}

\headsec{1. Узнайте повторяющуюся конструкцию}
Замена переменной полезна, когда одно и то же выражение встречается в уравнении несколько раз.
Вместо громоздкой конструкции \(u(x)\) вводят короткое обозначение \(t=u(x)\).
\[
  F\bigl(u(x)\bigr)=0\quad\longrightarrow\quad F(t)=0.
\]
\key{Обязательные шаги:} выбрать \(t\), записать его допустимые значения, решить уравнение для \(t\), выполнить обратную замену и проверить корни.

\headsub{Пример 1. Замена квадрата}
\[
 (x+2)^4-4(x+2)^2-5=0.
\]
Заметим, что \((x+2)^4=\bigl((x+2)^2\bigr)^2\). Положим
\[
 t=(x+2)^2,\qquad t\ge0.
\]
Тогда \(t^2-4t-5=0\), откуда \(t=5\) или \(t=-1\). Второе значение исключаем:
\(t\) обозначает квадрат действительного числа.
\[
 (x+2)^2=5\quad\Longleftrightarrow\quad x+2=\pm\sqrt5.
\]
\key{Ответ:} \(x=-2-\sqrt5\), \(x=-2+\sqrt5\).

\note{Важно: из \(A^2=c>0\) следуют оба равенства \(A=\sqrt c\) и \(A=-\sqrt c\).}

\headsub{Пример 2. Замена дроби целиком}
\[
 \frac{1}{(x-2)^2}-\frac{1}{x-2}-6=0,\qquad x\ne2.
\]
Пусть \(t=\dfrac1{x-2}\). Здесь \(t\ne0\), но отрицательные значения допустимы.
\[
 t^2-t-6=0\quad\Longleftrightarrow\quad(t-3)(t+2)=0.
\]
\[
 \begin{aligned}
 t=3&\Longrightarrow\frac1{x-2}=3\Longrightarrow x=\frac73;\\
 t=-2&\Longrightarrow\frac1{x-2}=-2\Longrightarrow x=\frac32.
 \end{aligned}
\]
\key{Ответ:} \(\dfrac32\), \(\dfrac73\).

\clearpage
\headsec{2. Проверьте: нужна ли замена вообще?}
Если одинаковые слагаемые стоят по разные стороны равенства, сначала посмотрите,
можно ли их сократить. \key{При этом ограничения исходного выражения сохраняются.}

\headsub{Пример 3. Радикалы исчезли, ОДЗ осталась}
\[
 x^2-2x+\sqrt{3-x}=8+\sqrt{3-x}.
\]
Из корня следует \(3-x\ge0\), то есть \(x\le3\). После сокращения одинаковых
радикалов получаем
\[
 x^2-2x-8=0\quad\Longleftrightarrow\quad(x-4)(x+2)=0.
\]
Кандидаты: \(4\) и \(-2\). При \(x=4\) радикал в исходном уравнении не определён,
поэтому этот кандидат исключается. \key{Ответ:} \(x=-2\).

\headsec{3. Узнавайте другие полезные структуры}
\headsub{Пример 4. Сумма двух неотрицательных выражений}
Для действительных \(A\) и \(B\):
\[
 A^2+B^2=0\quad\Longleftrightarrow\quad
 \begin{cases}A=0,\\B=0.\end{cases}
\]
Например,
\[
 (x^2-25)^2+(x^2+3x-10)^2=0
\]
превращается в систему
\[
 \begin{cases}x^2-25=0,\\x^2+3x-10=0.\end{cases}
 \qquad\Longleftrightarrow\qquad
 \begin{cases}x=\pm5,\\x=-5\ \text{или}\ x=2.\end{cases}
\]
Подходят только \emph{общие} решения обоих уравнений. \key{Ответ:} \(x=-5\).

\note{Тот же принцип работает для суммы модулей, а также для суммы арифметических корней,
если все подкоренные выражения неотрицательны.}

\headsub{Пример 5. Разность квадратов}
\[
 x^4=(4x-5)^2\quad\Longleftrightarrow\quad
 \underbrace{\bigl(x^2-(4x-5)\bigr)}_{A-B}
 \underbrace{\bigl(x^2+(4x-5)\bigr)}_{A+B}=0.
\]
Получаются два уравнения: \(x^2-4x+5=0\) (действительных корней нет) и
\(x^2+4x-5=0\) (корни \(-5\) и \(1\)).
\key{Ответ:} \(x=-5\), \(x=1\).

\note{Следите за скобками: \(x^2-(4x-5)=x^2-4x+5\).}

\clearpage
\headsec{4. Подготовьте выражение к замене}
Повторяющаяся часть иногда скрыта. Удачная перегруппировка позволяет получить одно и то же \(t\)
в нескольких местах уравнения.

\headsub{Пример 6. Разница в одну единицу}
\[
 (x^2-5x+7)^2-(x-2)(x-3)=3.
\]
Поскольку \((x-2)(x-3)=x^2-5x+6\), обозначим
\[
 t=x^2-5x+7=\left(x-\frac52\right)^2+\frac34,\qquad t\ge\frac34.
\]
Тогда \((x-2)(x-3)=t-1\) и
\[
 t^2-(t-1)=3\quad\Longleftrightarrow\quad
 t^2-t-2=0\quad\Longleftrightarrow\quad t=2\ \text{или}\ t=-1.
\]
Остаётся только \(t=2\). Обратная замена:
\[
 x^2-5x+7=2\quad\Longleftrightarrow\quad x^2-5x+5=0.
\]
\key{Ответ:} \(x=\dfrac{5-\sqrt5}{2}\), \(x=\dfrac{5+\sqrt5}{2}\).

\headsub{Пример 7. Похожие знаменатели}
\[
 \frac1{x(x+3)}-\frac1{x^2+3x+2}=\frac1{60}.
\]
Положим \(t=x^2+3x\). Тогда \(t\ne0\), \(t\ne-2\) и
\[
 \frac1t-\frac1{t+2}=\frac1{60}
 \quad\Longleftrightarrow\quad
 \frac2{t(t+2)}=\frac1{60}.
\]
Следовательно, \(t^2+2t-120=0\), то есть \(t=10\) либо \(t=-12\).
Однако
\[
 t=x^2+3x=\left(x+\frac32\right)^2-\frac94\ge-\frac94,
\]
поэтому \(t=-12\) невозможно. При \(t=10\):
\[
 x^2+3x-10=0\quad\Longleftrightarrow\quad x=2\ \text{или}\ x=-5.
\]
Оба значения допустимы в исходном уравнении. \key{Ответ:} \(x=-5\), \(x=2\).

\vspace{1mm}
\noindent\textbf{\color{darkaccent}Самопроверка.} Найдена удобная замена; сохранены ограничения;
учтены все ветви обратной замены; корни проверены в исходном уравнении.

\clearpage
\thispagestyle{plain}
\vspace*{9mm}
\begin{center}
{\Large\bfseries\color{darkaccent}Алгоритм решения уравнения заменой переменной}\par
\vspace{15mm}
\begin{tikzpicture}[node distance=12mm]
\node[fstart] (s) {Начало};
\node[faction,below=of s] (a) {Найдите повторяющуюся конструкцию \(u(x)\).\\При необходимости сгруппируйте или упростите выражения.};
\node[faction,below=of a] (b) {Введите \(t=u(x)\).\\Запишите ОДЗ исходного уравнения и ограничения для \(t\).};
\node[faction,below=of b] (c) {Получите уравнение относительно \(t\) и решите его.};
\node[faction,below=of c] (d) {Отберите допустимые значения \(t\).\\Для каждого из них решите \(u(x)=t\).};
\node[faction,below=of d] (e) {Проверьте найденные \(x\) в исходном уравнении.\\Исключите недопустимые значения.};
\node[fstart,below=of e] (f) {Запишите ответ};
\draw[farrow] (s)--(a);
\draw[farrow] (a)--(b);
\draw[farrow] (b)--(c);
\draw[farrow] (c)--(d);
\draw[farrow] (d)--(e);
\draw[farrow] (e)--(f);
\end{tikzpicture}
\end{center}
\clearpage

\headsec{Тренировочный блок · Путь А}
\note{Решите десять уравнений. Оформляйте замену, ограничения, обратную замену и ответ.
Задания расположены от простых к более содержательным.}
\vspace{4mm}
\begin{enumerate}[label=\textbf{\arabic*.},itemsep=7mm]
 \item \(\displaystyle (x+1)^4-5(x+1)^2+4=0.\)
 \item \(\displaystyle (x-3)^4-3(x-3)^2-4=0.\)
 \item \(\displaystyle \frac1{(x-2)^2}-\frac2{x-2}-3=0.\)
 \item \(\displaystyle \frac2{(x+1)^2}+\frac1{x+1}-3=0.\)
 \item \(\displaystyle x^2-3x+\sqrt{2-x}=4+\sqrt{2-x}.\)
 \item \(\displaystyle (x^2-9)^2+(x^2+x-12)^2=0.\)
 \item \(\displaystyle x^4=(2x+3)^2.\)
 \item \(\displaystyle (x^2-4x+4)^2-(x^2-4x+2)=4.\)
 \item \(\displaystyle x^2-2x+2=\frac2{x^2-2x+1}.\)
 \item \(\displaystyle \frac1{x(x+4)}-\frac1{x^2+4x+3}=\frac1{18}.\)
\end{enumerate}
\vfill
\note{Упражнения нацелены на распознавание структуры: прежде чем раскрывать скобки или
приводить дроби к общему знаменателю, попробуйте найти удачную замену.}
\clearpage

\headsec{Ответы к тренировочному блоку}
\renewcommand{\arraystretch}{1.38}
\begin{tabularx}{\linewidth}{@{}c>{\raggedright\arraybackslash}X>{\raggedright\arraybackslash}p{0.35\linewidth}@{}}
\toprule
\textbf{№} & \textbf{Ответ} & \textbf{Ключ к самопроверке}\\
\midrule
1 & \(x=-3,-2,0,1\) & \(t=(x+1)^2\ge0\)\\
2 & \(x=1,5\) & \(t=(x-3)^2\ge0\)\\
3 & \(x=1,\;\dfrac73\) & \(x\ne2\)\\
4 & \(x=-\dfrac53,\;0\) & \(x\ne-1\)\\
5 & \(x=-1\) & Учитывайте \(x\le2\)\\
6 & \(x=3\) & Требуется пересечение корней двух уравнений\\
7 & \(x=-1,3\) & Разность квадратов\\
8 & \(x=2-\sqrt2,\;2+\sqrt2\) & \(t=(x-2)^2\ge0\)\\
9 & \(x=0,2\) & \(t=(x-1)^2>0\)\\
10 & \(x=-2-\sqrt{10},\;-2+\sqrt{10}\) & \(t=x^2+4x\ge-4\)\\
\bottomrule
\end{tabularx}
\vspace{10mm}
\textbf{Особенно проверьте:}
\begin{itemize}
\item В задаче 5 корень \(x=4\) возникает после сокращения радикалов, но нарушает ОДЗ.
\item В задачах 8--10 ограничения для \(t\) помогают исключать лишние ветви.
\item В задачах 3, 4, 9 и 10 проверка знаменателей обязательна.
\end{itemize}
\vfill
\begin{center}\small\color{quiet}
Лёвин Артём Александрович эксклюзивно для Ярослава Верещагина
\end{center}
\end{document}
